\section*{Two operators with the same scalar shadow}
\addcontentsline{toc}{section}{Two operators with the same scalar shadow}

Return to uniform three-state references and consider two gauged blocks
\begin{equation}
G_A=
\begin{pmatrix}
2&-1&-1\\
-1&1&0\\
-1&0&1
\end{pmatrix},
\qquad
G_B=\sqrt{\frac53}
\begin{pmatrix}
0&1&-1\\
1&-1&0\\
-1&0&1
\end{pmatrix}.
\label{eq:book-same-shadow-blocks}
\end{equation}
Both have zero row and column means.  Their ordinary squared Frobenius norms are both $10$, so
their weighted scalar shadows are identical:
\begin{equation}
M_A=M_B=\frac{\sqrt{10}}{3}.
\label{eq:book-same-shadow-norm}
\end{equation}

The operators are nevertheless different.  The two nonzero singular values of the whitened
$G_A$ are $1$ and $1/3$, whereas those of the whitened $G_B$ are both $\sqrt5/3$.  One concentrates
more strength in a leading substitution direction; the other distributes its strength equally
across the two-dimensional contrast fiber.

Their action on the same source contrast $f=(1,-1,0)^{\mathsf T}$ makes the difference explicit:
\begin{equation}
\mathcal K_Af
=
\begin{pmatrix}
1\\[-2pt]-2/3\\[-2pt]-1/3
\end{pmatrix},
\qquad
\mathcal K_Bf
=\frac{1}{3}\sqrt{\frac53}
\begin{pmatrix}
-1\\[-2pt]2\\[-2pt]-1
\end{pmatrix}.
\label{eq:book-same-shadow-actions}
\end{equation}
The first operator favors the first target contrast; the second suppresses it and favors the
second.  A contact map built only from $M_{ij}$ assigns the two edges exactly the same value.

This is not an exotic collision.  Sign reversal, rotation of singular vectors, and redistribution
among singular values all create further operators with the same norm.  Scalarization is suitable
for asking where strong pair structure is located.  It cannot answer which substitutions are
compatible, compensatory, or antagonistic.  Those are operator-level questions.
